\begin{notechapter}
\chapter{Information Geometry and the Natural Gradient}
\label{ch:natural-gradient}
\chaptersubtitle{%
  The metric that a statistical model brings with it}

\begin{chapterplan}
\begin{itemize}[nosep]
  \item Statistical models as manifolds; the Fisher information as a Riemannian
        metric; its invariance under reparametrization. (Expectations are used
        here in the elementary sense and made rigorous in the next part.)
  \item The Kullback--Leibler divergence and its second-order expansion.
  \item Natural gradient descent; exponential families; the link with
        Gauss--Newton.
  \item Mirror descent as a second geometry-aware method.
\end{itemize}
\smallskip\noindent\textbf{Machine learning threads.}
\begin{itemize}[nosep]
  \item Kronecker-factored approximations of the Fisher matrix.
  \item Adam as a diagonal approximation, and what that costs in invariance.
  \item Trust regions in policy optimization.
\end{itemize}
\end{chapterplan}
\end{notechapter}
